\documentclass[a4paper, serif, 10pt]{moderncv}

%% moderncv
% color and style
\moderncvcolor{black}
\moderncvstyle{banking}

%% page layout/margins
\usepackage[top=2.5cm, bottom=2.5cm, left=1.5cm, right=1.5cm]{geometry}

\nopagenumbers{}

%% Babel config for English language
\usepackage[english]{babel}

%% fontspec
\usepackage{fontspec}

\IfFontExistsTF{Linux Libertine}{
  \setmainfont[Ligatures={TeX, Common}, Numbers=OldStyle]{Linux Libertine}
}{}
\IfFontExistsTF{Linux Libertine O}{
  \setmainfont[Ligatures={TeX, Common}, Numbers=OldStyle]{Linux Libertine O}
}{}

%% CTeX
% CJK support, used to show CJK characters in the resume
%
% - heading=false: disable ctex heading
% - punct=kaiming: use kaiming punctuations styles for CJK
% - scheme=plain: use plain scheme, do not override \normalsize font size
% - space=auto: space settings for CJK characters
%
% CJK fonts are resolved per locale language, e.g. Noto Serif/Sans CJK JP
% for Japanese resumes, with fallback to the next region if not installed
%
% ref:
% - http://ctan.mirrorcatalogs.com/language/chinese/ctex/ctex.pdf
\usepackage[UTF8, heading=false, punct=kaiming, scheme=plain, space=auto]{ctex}

\IfFontExistsTF{Noto Serif CJK SC}{
  \setCJKmainfont{Noto Serif CJK SC}
}{}
\IfFontExistsTF{Noto Sans CJK SC}{
  \setCJKsansfont{Noto Sans CJK SC}
}{}

%% line spacing
\usepackage{setspace}
\setstretch{1.125}

%% URL styling - use normal text instead of monospace
\urlstyle{same}

%% Patch \httplink and \httpslink to handle full URLs
% The original moderncv macros always prepend http:// or https:// to URLs.
% This causes issues when the URL already has a protocol (e.g., https://example.com)
% resulting in malformed URLs like https://https://example.com.
% This patch checks if the URL already contains "://" and uses it directly if so.
% Note: We use \str_set:Nx (x-expansion) to expand macros like \@homepage first.
\ExplSyntaxOn
\str_new:N \g_yamlresume_url_str

\RenewDocumentCommand{\httpslink}{O{}m}{
  \str_set:Nx \g_yamlresume_url_str {#2}
  \str_if_in:NnTF \g_yamlresume_url_str {://}
    { \url{#2} }
    { \url{https://#2} }
}

\RenewDocumentCommand{\httplink}{O{}m}{
  \str_set:Nx \g_yamlresume_url_str {#2}
  \str_if_in:NnTF \g_yamlresume_url_str {://}
    { \url{#2} }
    { \url{http://#2} }
}
\ExplSyntaxOff

\name{陳嘉俊}{}
\title{人工智能工程師 | 生成式 AI 與機器學習}
\phone[mobile]{+852 9123 4567}
\email{ka.chun.chan@example.com}
\homepage{https://www.kachun-ai.dev}

\address{香港中環皇后大道中 99 號, 香港, 香港島, Hong Kong}{}{}

\extrainfo{{\small \faGithub}\ \href{https://github.com/kachun-ai}{@kachun-ai} {} {} {} • {} {} {} 
{\small \faLinkedin}\ \href{https://www.linkedin.com/in/kachun-ai}{@kachun-ai} {} {} {} • {} {} {} 
{\small \faMedium}\ \href{https://medium.com/@kachun-ai}{@kachun-ai}}

\begin{document}

\maketitle

\section{Basics}

\cvline{}{\begin{itemize}
\item 人工智能工程師，專注生成式 AI、大型語言模型與檢索增強生成（RAG）應用落地
\item 具備 5 年以上機器學習與數據工程經驗，熟悉模型訓練、評估、部署及 MLOps 全流程
\item 擅長以 Python、PyTorch 及雲端平台建構可擴展的 AI 服務，並推動負責任 AI 實踐
\item 具備跨部門協作與技術指導經驗，能將業務需求轉化為可量測的 AI 解決方案
\end{itemize}}
\section{Education}

\cventry{Sep 2018–Jul 2020}
        {Master, 計算機科學與工程, Score: 3.9/4.0}
        {香港科技大學}
        {\url{https://www.ust.hk}}
        {}
        {\begin{itemize}
\item 主修機器學習與自然語言處理，深入研究大型語言模型與檢索增強生成技術
\item 以優異成績完成畢業論文，主題為粵語語音辨識與意圖理解
\item 積極參與人工智能實驗室研究項目，累積模型訓練與部署經驗
\end{itemize}
\textbf{Courses}: 機器學習, 深度學習, 自然語言處理, 分散式系統, 雲端計算, 資料探勘, 演算法設計與分析}

\cventry{Sep 2014–Jun 2018}
        {Bachelor, 計算機科學, Score: 3.7/4.0}
        {香港中文大學}
        {\url{https://www.cuhk.edu.hk}}
        {}
        {\begin{itemize}
\item 建立扎實的計算機科學基礎，熟悉演算法、資料結構與軟件工程
\item 曾參與校內人工智能研究小組，開發影像分類與文字分析專案
\item 獲列入學院院長榮譽名單
\end{itemize}
\textbf{Courses}: 資料結構, 作業系統, 資料庫系統, 機器學習導論, 統計學, 線性代數}
\section{Work}

\cventry{Mar 2022–Present}
        {高級人工智能工程師}
        {智言科技有限公司}
        {\url{https://www.zhiyan-ai.com}}
        {}
        {\begin{itemize}
\item 帶領 4 人團隊設計與落地企業級檢索增強生成（RAG）平台，支援粵語、英語及普通話問答
\item 建立大型語言模型評估流程，涵蓋準確度、幻覺率、延遲與成本指標，將模型上線週期縮短 40\%
\item 開發基於 Python、PyTorch 與 LangChain 的模型服務，並以 Docker 與 Kubernetes 部署至雲端
\item 與產品、法務及合規團隊協作，制定負責任 AI 與資料私隱審查規範
\item 指導初級工程師與數據科學家，推動程式碼審查、實驗追蹤與 MLOps 最佳實踐
\end{itemize}
\textbf{Keywords}: 大型語言模型, 檢索增強生成, MLOps, Python, 模型評估, Kubernetes}

\cventry{Jul 2020–Feb 2022}
        {機器學習工程師}
        {香港數碼智能有限公司}
        {\url{https://www.hkdigitalai.com}}
        {}
        {\begin{itemize}
\item 開發客戶流失預測與推薦系統，運用 XGBoost、LightGBM 及深度學習模型提升轉換率 18\%
\item 建置自動化特徵工程與模型訓練管道，使用 Airflow、Spark 及 MLflow 管理實驗與版本
\item 將模型封裝為 RESTful API，部署於 AWS 並整合實時資料串流，支援每日百萬級預測請求
\item 與業務分析師合作定義關鍵指標，透過 A/B 測試驗證模型成效並持續迭代
\end{itemize}
\textbf{Keywords}: 機器學習, XGBoost, Spark, MLflow, AWS, A/B 測試}
\section{Skills}

\cvline{人工智能與機器學習}{Expert \hfill \textbf{Keywords}: 深度學習, 自然語言處理, 大型語言模型, 檢索增強生成, 模型評估, 電腦視覺}
\cvline{程式語言與框架}{Advanced \hfill \textbf{Keywords}: Python, PyTorch, TensorFlow, scikit-learn, LangChain, SQL}
\cvline{資料工程與 MLOps}{Advanced \hfill \textbf{Keywords}: Docker, Kubernetes, MLflow, Airflow, Spark, AWS, GCP}
\end{document}